% ===================== 第 4 章 =====================
\section{七人科学议会}

\subsection{七个固定席位}

每个正式研究任务中，七个席位全程参与，定义在 \code{packages/council/contracts.py} 的 \code{Seat} 枚举。席位不是可选项标签，而是各自拥有独立角色规格（\code{roles.py::ROLE\_SPECS}）、独立私人记忆（agent id 为 \code{\{task\_id\}:\{seat\}}，没有任何代码读取他人 id）和不同证据排序权重（\code{packages/memory/projection.py}）。这三项差异都有测试锁定——历史上曾出现五个席位共用一份空权重、在唯一用于区分它们的维度上互为复制品的事故。

\begin{table}[ht]
\centering\small
\caption{七个席位与其质询重心}
\label{tab:seats}
\begin{tabularx}{\linewidth}{@{}L{4.0cm} X@{}}
\toprule
\rowcolor{accent}
\tblhead 席位（标识） & \tblhead 质询重心 \\
\midrule
理论建构者 \seat{theory\_builder}            & 竞争理论、机制链条、风险预测是否成立 \\
因果推断专家 \seat{causal\_scientist}        & 混杂、反向因果、选择偏差、识别策略 \\
测量与构念专家 \seat{measurement\_scientist} & 构念与操作化的距离、自报告偏差、测量不变性 \\
统计与复现专家 \seat{replication\_scientist} & 功效、精度、多重比较、复现史 \\
边界与情境专家 \seat{boundary\_scientist}    & 适用范围、异质性、外推条件 \\
对抗性证伪者 \seat{adversarial\_falsifier}   & 攻击结论的最强版本，主动寻找反例与零结果 \\
证据与溯源审计员 \seat{evidence\_auditor}    & 来源锚点、引用蕴含、证据独立性 \\
\bottomrule
\end{tabularx}
\end{table}

议会之上没有第八名科学家：调度、记录与组装由编排层承担，它\textbf{无投票权}，不代表任何学术立场；\code{EpistemoBrain} 则位于编排层之下的记忆层，只负责记忆的压缩与重建，同样不参与任何裁决。七席共享同一套运行框架（\code{deliberation.py} 中唯一的 \code{GatewayDeliberator} 类服务全部席位）与同一个工具缓存，既满足「不复制七套业务代码、不部署七个服务」，又保住席位之间的独立性。

\subsection{八阶段协议}

阶段顺序由 \code{packages/epistemo/contracts.py} 的 \code{PHASE\_SEQUENCE} 冻结，状态机只允许顺序前进（第 9.1 节）。

\begin{enumerate}[leftmargin=1.6em, itemsep=3pt]
\item \textbf{独立预承诺。}七席在看到他人意见前，分别提交初始判断、置信度、三个最可能的盲点、待补证据与改判条件。预承诺封存后才允许读取，读取后不能补交——后续一致性因此无法来自锚定。
\item \textbf{专业取证。}各席按自身维度提出证据需求，查询规划器合并去重后批量检索（第 8.4 节），驱动搜索的是证据需求矩阵而不是泛泛关键词。
\item \textbf{证据交换。}交换结构化证据（主张、原文片段、研究设计、质量等级、适用边界），不交换完整思维链。
\item \textbf{交叉质询。}目标是暴露隐藏假设而非争胜。被质询方只能回应 \code{DEFEND / REVISE / NARROW / WITHDRAW / DISSENT} 五种动作之一。
\item \textbf{盲点悬赏。}证据图暴露的缺口生成盲点候选，按「影响 $\times$ 不确定性 $\times$ 可调查性 $\times$ 新颖度 $-$ 成本」打分，广播给七席共同认领。
\item \textbf{联合建模。}把各席结论组合为条件化共识；最强反对或证伪条件缺失时拒绝出共识。进入本阶段前有唯一的人类引导检查点（第 4.5 节）。
\item \textbf{最终复判。}七席再次独立判断，记录信念更新轨迹；仍持异议的席位获得具名异议证书。复判只遍历实际参与的席位，不为缺席者伪造判断（这条由评测基线的一次事故促成，见第 10.5 节）。
\item \textbf{报告组装。}从图与账本组装产物，本身不产生新的模型判断。
\end{enumerate}

一个席位失败不会拖垮任务：失败被记为该轮 \code{SEAT\_UNAVAILABLE} 事件与一个未填槽位，其余席位继续，任务最终以 \code{COMPLETED\_WITH\_GAPS} 收尾，缺席者在报告局限一节逐条点名（第 9.4 节）。

\subsection{结构化科研动作：科学家能说什么是被限定的}

科学家不能向账本写自由发言，只能提交八类结构化动作；收到质询后，回应动作再收窄为五类。

\begin{center}\small
\begin{tcolorbox}[colback=panel, colframe=linegray, boxrule=0.6pt, arc=2pt, left=8pt,right=8pt,top=6pt,bottom=6pt, fontupper=\ttfamily]
主动作：PROPOSE\quad SUPPORT\quad CHALLENGE\quad QUALIFY\quad FORK\quad REQUEST\quad REVISE\quad DISSENT\\[3pt]
应答：DEFEND\quad REVISE\quad NARROW\quad WITHDRAW\quad DISSENT
\end{tcolorbox}
\end{center}

每个动作必须携带 actor、target、statement、\code{evidence\_refs}、confidence、\code{falsification\_condition}、novelty。校验器（\code{council/actions.py}）会拒绝：没有目标节点的泛泛评论、没有来源的事实主张、与既有内容高度重复的观点、只有「我同意」而无新增信息的发言、以及说不出如何被推翻的绝对断言。动作词汇表本身就是一份「科研讨论该长什么样」的形式化约束。

\subsection{禁止多数投票：共识如何被阻断}

在 \code{evaluate\_consensus} 中，即使多数席位投出 \code{SUPPORT}，只要存在未解决的致命质询（\code{has\_unresolved\_fatal\_challenge}）或证据引用为空，共识状态就是 \code{BLOCKED}。联合建模处理器在最强反对与证伪条件缺失时同样拒绝产出共识。因此最终输出不是票数，而是四段式结构：

\begin{mechbox}
\textbf{多数判断}：存在小幅、条件性的影响。\\
\textbf{少数异议}：现有研究不足以排除反向因果。\\
\textbf{真正分歧}：纵向时间顺序能否被视为充分的因果证据。\\
\textbf{解决分歧所需证据}：带平台日志与基线心理健康控制的预注册纵向研究。
\end{mechbox}

\subsection{人类方向性引导检查点}

盲点悬赏结束、联合建模开始之前，任务进入 \code{AWAITING\_COUNCIL\_INPUT}：研究者可以查看七席此刻的立场（预承诺、置信度、已提质询），提交一段方向性备注，也可以明确留空。这个检查点有四条边界，保证它不违反「禁止投票」：它是单一研究者的调度性输入而非票数；只影响下一轮重点讨论哪些冲突；不改变任何证据门判定；不会作为证据进入任何正式节点。备注在联合建模提示词中以独立、显式标注的文本块注入（如「[研究者方向性备注，非科学判断]：……」）。

研究者离开页面不会卡住任务：宽限期默认 900 秒（15 分钟；环境变量全名见附录 A.4，下限 30 秒），届满自动以「无干预」回到队列续跑。实现横跨状态枚举、\code{CouncilCheckpoint.guidance} 字段、预览与引导两个 API 端点（\code{council-preview}、\code{council-guidance}）、CLI 子命令与前端 \code{CheckpointGate.tsx}。
